\section{Introduzione}

\subsection{Scopo del documento}
Lo scopo del documento è descrivere l'analisi dei requisiti del progetto \textbf{Second Brain}. Il testo definisce in modo strutturato i casi d’uso e le funzionalità richieste, classificandole in requisiti obbligatori, desiderabili e opzionali, secondo le specifiche fornite dal capitolato d'appalto C6 proposto dall'azienda Zucchetti.

\subsection{Scopo del prodotto}
Il prodotto mira a fornire un editor Markdown intelligente che funga da vero e proprio "secondo cervello" per l'utente, aiutandolo a organizzare e valorizzare i propri spunti creativi e appunti estemporanei. L'applicazione estende le normali capacità di note-taking integrando la potenza dei Large Language Models (LLM). Oltre a supportare la scrittura e la formattazione grafica, il sistema funge da assistente attivo per il brainstorming, l'analisi critica dei testi tramite il metodo dei "Sei Cappelli per Pensare" di Edward De Bono, e la generazione automatica di contenuti (Distant Writing). Il prodotto finale sarà uno strumento accessibile anche a utenti non esperti.

\subsection{Glossario}
Al fine di evitare ambiguità, i termini tecnici o con significato specifico sono seguiti dal pedice e definiti nel documento \link{https://synergo-unit-unipd.github.io/Documentary-Repo/docs/RTB/doc_gestionali/doc_interni/glossario/glossario.pdf}{Glossario}.

\subsection{Riferimenti}
\subsubsection{Normativi}
\begin{itemize}
	\item \link{https://www.math.unipd.it/~tullio/IS-1/2025/Progetto/C6.pdf}{Capitolato C6 - Progetto Second Brain}
	\item \link{https://www.math.unipd.it/~tullio/IS-1/2025/Dispense/PD1.pdf}{Regolamento progetto didattico}
	\item \link{https://synergo-unit-unipd.github.io/Documentary-Repo/docs/RTB/doc_gestionali/doc_interni/norme_di_progetto/norme_di_progetto.pdf}{Norme di Progetto}
\end{itemize}

\subsubsection{Informativi}
\begin{itemize}
	\item Slide del corso di Ingegneria del Software (T05 - Analisi dei Requisiti)
	\item Verbali esterni ed interni
	\item E. De Bono, \textit{Sei cappelli per pensare} (riferimento metodologico per l'analisi critica)
    \item Documentazione API OpenAI (standard di comunicazione LLM richiesto dal proponente)
\end{itemize}